% Generated by docs/user-manual/tables/model_coverage.py -- do not hand-edit.
% Source: day 0 of releases 2026-11-06 to 2026-11-13.
\begin{table}[htbp]
\centering
\small
\setlength{\tabcolsep}{3pt}
\begin{tabular}{l r r r r r r}
\toprule
River model & Share of alerts & Warning & Danger & Extreme & 2-model & 3-model \\
\midrule
  GEOGLOWS & 8.5\% -- 27.2\% & 79.4\% -- 86.9\% & 7.4\% -- 14.9\% & 5.1\% -- 7.6\% & 0.6\% -- 4.5\% & 0.0\% \\
  Flood Hub & 0.2\% -- 0.4\% & 33.3\% -- 100.0\% & 0.0\% -- 66.7\% & 0.0\% & 0.0\% -- 12.5\% & 0.0\% \\
  GloFAS & 72.7\% -- 91.4\% & 68.9\% -- 80.1\% & 11.6\% -- 19.4\% & 4.4\% -- 15.1\% & 0.1\% -- 1.5\% & 0.0\% \\
\bottomrule
\end{tabular}
\caption[What each river model contributes]{%
  \textbf{What each river model contributes.}
  Every figure is the low--high range across day~0 of
  8 releases (2026-11-06 to 2026-11-13).
  \emph{Share of alerts} is each model's portion of all river forecasts
  issued for that day, so the column adds to 100\%. The severity columns are
  the mix \emph{within} each model, so each row adds to 100\%. Shares say
  who issues the alerts, not who is right.
  The last two columns change denominator: they are the share of that model's
  own flagged \emph{cells} that also carry one, or two, other models --
  corroboration is spatial, so it is counted on cells rather than on
  alerts. GEOGLOWS figures exclude the blacked-out region in 7 of 8 releases (2026-11-06 to 2026-11-12): 427,826 forecast(s) withheld in total.}
\label{tab:model-coverage}
\end{table}
